\chapter{Introduction}

    \paragraph{}
    In an engineering environment as intricate and aspirational as Sunswift Racing, where hundreds of thousands of requirements intersect with evolving design decisions, regulatory obligations and subsystem interactions, our research is guided by a single unifying question that animates every stage of this thesis: How can an ontology driven, knowledge graph enabled, large language model supported framework be designed to transform fragmented regulatory and engineering artefacts into a coherent, explainable and adaptive system for regulatory aligned automotive design? To answer this question we situate our work at the meeting point between structured semantic representation and expressive neural reasoning, allowing us to explore how formal ontologies, knowledge graphs and large language models may collaborate to illuminate obligations hidden within dense regulatory text, uncover dependencies woven through complex vehicle architectures and ultimately support engineers in designing systems that are both compliant and conceptually transparent. The thesis therefore proceeds not merely as an exercise in technical construction but as an attempt to breathe interpretative life into the sprawling documentary landscape of the Australian Design Rules, transforming them from static, isolated documents into a living network of meaning, inference and design intelligence.

    \paragraph{}
    Because this research spans the full arc of knowledge creation from source understanding to context aware inference, we divide our explorations between two complementary streams that reflect the dual nature of our methodological ambition: Manan focuses on knowledge extraction, designing the ontology driven systems engineering framework that anchors the entire structure, developing the model context foundations that integrate structured (knowledge graph) and unstructured (language) sources, and establishing the mechanisms through which risk, regulatory alignment and cross subsystem traceability can be formalised with precision; Ilina focuses on knowledge extension, constructing the pipeline through which large language models, neural link prediction and reinforcement learning agents expand the graph beyond the limits of literal text, transforming legacy documentation into an intelligent, queryable knowledge base capable of offering explainable compliance, risk and design recommendations. Together these two streams form a single coherent endeavour, in which extraction provides the discipline and structural clarity upon which the extension processes rely, and extension offers the interpretative breadth, adaptivity and inferential grace that bring the extracted knowledge into meaningful engagement with Sunswift’s real engineering needs.